\documentclass[11pt,a4paper]{article}

% ============================================================
% Packages (core only --- texlive-latex-base + texlive-base)
% ============================================================
\usepackage{amsmath, amssymb, amsthm}
\usepackage[a4paper, margin=2.5cm]{geometry}
\usepackage{array}

% ============================================================
% Theorem environments
% ============================================================
\theoremstyle{definition}
\newtheorem{definition}{Definition}[section]
\newtheorem{remark}{Remark}[section]
\newtheorem{example}{Example}[section]

\theoremstyle{plain}
\newtheorem{proposition}{Proposition}[section]
\newtheorem{lemma}{Lemma}[section]
\newtheorem{corollary}{Corollary}[section]

% ============================================================
% Notation macros
% ============================================================
\newcommand{\R}{\mathbb{R}}
\newcommand{\E}{\mathbb{E}}
\newcommand{\N}{\mathbb{N}}
\newcommand{\norm}[1]{\left\lVert #1 \right\rVert}
\newcommand{\inner}[2]{\langle #1,\, #2 \rangle}
\newcommand{\corr}{\operatorname{corr}}
\newcommand{\cov}{\operatorname{cov}}
\newcommand{\Var}{\operatorname{Var}}
\newcommand{\med}{\operatorname{median}}
\newcommand{\sgn}{\operatorname{sign}}
\newcommand{\Neff}{N_{\mathrm{eff}}}
\newcommand{\ICbar}{\overline{\mathrm{IC}}}
\newcommand{\rhobar}{\bar{\rho}}
\newcommand{\given}{\,\vert\,}
\DeclareMathOperator{\MI}{I}
\DeclareMathOperator{\TE}{TE}
\DeclareMathOperator{\IR}{IR}
\DeclareMathOperator{\IC}{IC}
\DeclareMathOperator{\BR}{BR}
\DeclareMathOperator*{\argmax}{arg\,max}
\DeclareMathOperator*{\argmin}{arg\,min}
\newcommand{\code}[1]{\texttt{\detokenize{#1}}}

% ============================================================
\begin{document}

\title{\textbf{The Process: A First-Class Analytical Unit for\\[3pt]
Orthogonal Information Extraction from Market Microstructure}\\[6pt]
\large Conditional-Information Signal Selection and the Effective\\
Degrees of Freedom of an Alpha Book}

\author{Yigit Onat\\
\textit{Independent Researcher}\\
\texttt{yionat@gmail.com}}
\date{June 2026}
\maketitle

% ============================================================
\begin{abstract}
A high-frequency microstructure feed produces hundreds of engineered features, most of them
mutually redundant. The operative research question is therefore not whether an individual feature
predicts future price action, but how much \emph{non-redundant} information the feature set carries
in aggregate, and how many statistically \emph{independent} bets that information supports. We
formalize this question around a third first-class analytical unit---the \emph{process}---placed
beside the \emph{feature} (what is computed) and the \emph{algorithm} (how it trades). A process is
a typed, registry-bound analytical contract that certifies whether a feature, or a
derivative/combination of features, carries information about a precisely specified target. We then
develop the mathematics the unit exists to serve. Adopting mutual information as the currency of
prediction (justified by Fano's inequality, which makes information a hard lower bound on
achievable error), we define the effective degrees of freedom of a signal book through the
fundamental law of active management, $\IR \approx \ICbar\,\sqrt{\Neff}$ with
$\Neff = N/[1+(N-1)\rhobar]$, and show that the binding quantity is the \emph{stress-regime}
correlation, not the calm-regime correlation. We give five orthogonalization mechanisms, each as a
process: conditional-mutual-information greedy selection, interaction-information synergy mining,
training-prefix residualization, Marchenko--Pastur-denoised principal components, and a capstone
decorrelated signal book that reports its own honest $\Neff$. We make ``price action'' explicit
through an aspect$\times$horizon target taxonomy, and we add a derivative-operator layer that mints
new feature degrees of freedom (fractional differentiation, leaky integration, causal spectral
features) under a program-level false-discovery ledger. The contribution is a unifying analytical
contract together with the orthogonality-and-breadth mathematics for disciplined microstructure
alpha discovery. This is a methods/framework paper: empirical estimates are specified as future
work, and every unit ships behind a mandatory planted-signal conformance gate.
\end{abstract}

\smallskip
\noindent\textbf{Keywords:} market microstructure, mutual information, conditional mutual
information, feature selection, effective breadth, fundamental law of active management, random
matrix theory, false discovery rate, fractional differentiation.

\smallskip
\noindent\textbf{JEL Classification:} C14, C52, C58, G14, G17

\tableofcontents
\bigskip

% ============================================================
\section{Introduction}
\label{sec:intro}
% ============================================================

\subsection{Redundancy is the problem, not scarcity}

A modern microstructure pipeline does not suffer from a shortage of candidate predictors. The
platform underlying this work emits $236$ engineered features every $100$\,milliseconds---order-book
imbalance at several depths, trade-flow toxicity, realized-volatility estimators, entropy of book
shape, funding pressure, liquidation geometry, and dozens of derived interactions. The difficulty
is the opposite of scarcity: these features are heavily \emph{co-determined}. Imbalance at level one
and level two move together; a dozen volatility proxies share a common factor; flow and price are
mechanically linked. Screening each feature in isolation and keeping ``the ones with a high
information coefficient'' therefore answers the wrong question. It rewards a crowd of near-copies of
a single underlying bet and quietly overstates how much independent evidence the research has
actually accumulated.

The correct question has two parts, both of which this paper formalizes:

\begin{enumerate}
  \item \textbf{Orthogonal extraction.} From a large, redundant feature set, isolate a basis of
        signals each of which carries information about future price action that the others do
        \emph{not}. Redundancy must be measured in an information metric, not a linear one, because
        microstructure dependence is routinely nonlinear (a feature can be linearly uncorrelated
        with the target yet fully determine it---the canonical XOR obstruction).
  \item \textbf{Effective degrees of freedom.} Quantify how many statistically independent bets the
        resulting book supports. This number---not the raw count of ``significant'' features---is
        what governs the achievable information ratio, and it is the quantity that collapses in
        exactly the market states where a portfolio most needs diversification.
\end{enumerate}

\subsection{Three first-class citizens}

The platform answers these questions by promoting the analytical act itself to a typed,
versioned, testable object. There are three first-class citizens, summarized in
Table~\ref{tab:citizens}. The \emph{feature} and the \emph{algorithm} are familiar; the
\emph{process} is the unit this paper defines.

\begin{table}[h]
\centering
\renewcommand{\arraystretch}{1.3}
\setlength{\tabcolsep}{6pt}
\begin{tabular}{>{\raggedright\arraybackslash}p{2.4cm}
                >{\raggedright\arraybackslash}p{6.4cm}
                >{\raggedright\arraybackslash}p{5.6cm}}
\hline
\textbf{Citizen} & \textbf{Question it answers} & \textbf{Contract} \\
\hline
\textbf{feature}   & \emph{what} is computed                                   & ingestor columns / derived series \\
\textbf{algorithm} & \emph{how} it trades                                      & \code{MicrostructureAlgorithm} \\
\textbf{process}   & \emph{whether/where} information about price action exists & \code{Process} \\
\hline
\end{tabular}
\caption{The three first-class citizens. A process formalizes what previously lived as scattered
one-off analyses (a screener's information coefficient, an information-theory engine's mutual
information, a spectral persistence study, a gradient-boosted importance test) behind one registry,
one result schema, one command-line surface, and one persistence path.}
\label{tab:citizens}
\end{table}

\subsection{Contributions}

\begin{enumerate}
  \item A formal, registry-bound \emph{process contract} (Section~\ref{sec:process}): two kinds
        (evaluation and transform), a versioned result schema, and a typed-node refinement that
        treats the prediction target as its own first-class object.
  \item The \emph{information currency} (Section~\ref{sec:info}): why mutual information, not
        correlation, is the right dependence measure, anchored in Fano's inequality and a
        cost-aware minimum-information gate.
  \item The \emph{effective degrees of freedom} of a signal book (Section~\ref{sec:dof}): a
        derivation of $\Neff = N/[1+(N-1)\rhobar]$ from the fundamental law of active management,
        and the calm-versus-stress decomposition that turns it into a risk disclosure.
  \item Five \emph{orthogonalization processes} (Section~\ref{sec:ortho}), each defined
        mathematically: conditional-MI greedy selection, interaction-information synergy mining,
        residualization, Marchenko--Pastur-denoised PCA, and the capstone signal book.
  \item A \emph{target taxonomy} (Section~\ref{sec:target}) making ``price action'' explicit along
        an aspect$\times$horizon grid, and a \emph{derivative-operator} layer
        (Section~\ref{sec:ops}) that mints new feature degrees of freedom under a program-level
        false-discovery ledger.
  \item An \emph{estimation-hygiene} layer and the planted-signal verification discipline
        (Sections~\ref{sec:hygiene}--\ref{sec:status}) that keep the degrees-of-freedom estimate
        honest.
\end{enumerate}

This is deliberately a methods/framework paper. Where the contract is implemented we say so; where a
mechanism is specified but not yet built we say so; and we report no fabricated empirical estimates.
The one empirical fact we import as motivation---that a mid-price information coefficient of
$\approx 0.45$ collapses to $\approx 0.03$ once conditioned on states where a passive fill would
actually occur---is established in a companion study and cited, not re-derived here.

% ============================================================
\section{The Process as an Analytical Unit}
\label{sec:process}
% ============================================================

\subsection{Definition and the two kinds}

\begin{definition}[Process]
\label{def:process}
A \emph{process} is an analytical map---statistical, signal-processing, or machine-learning---that
takes a panel of feature columns and a context, and certifies whether a feature, or a
derivative/combination of features, carries information about a specified target of future price
action. It is identified by a unique name, registered through a decorator, parameterized through a
configuration section keyed by that name, and persisted through a single result schema.
\end{definition}

Two kinds share the registry, command-line, and persistence surface.

\begin{definition}[Evaluation and transform processes]
An \emph{evaluation process} is a map
\[
   \mathcal{E} : (\,\mathbf{D},\, \mathcal{C}\,) \longmapsto \mathcal{R},
\]
where $\mathbf{D}\in\R^{T\times p}$ is a time-indexed feature panel, $\mathcal{C}$ is a
\emph{process context} (symbol, timeframe, price column, horizon set, cost parameters, target,
auxiliary sources), and $\mathcal{R}$ is a \emph{process result}: it scores existing columns and
produces no new series. A \emph{transform process} is a map
\[
   \mathcal{T} : (\,\mathbf{D},\, \mathcal{C}\,) \longmapsto (\,\mathbf{D}',\, \mathcal{R}\,),
   \qquad \mathbf{D}' \in \R^{T\times m},
\]
producing $m$ \emph{new} series $\mathbf{D}'$ on the input index (e.g.\ principal components,
triple-barrier labels, residualized innovations). Derived series are themselves first-class
evaluable inputs: the runner chains $\mathbf{D}'$ into any evaluation process, so
$\mathcal{E}\circ\mathcal{T}$ is itself a process.
\end{definition}

The closure of evaluation under transform is the engine of the whole framework: orthogonalization
is a transform whose output is then scored by an evaluation process, so ``extract the orthogonal
part, then ask what it knows'' is a single composed pipeline rather than two disconnected scripts.

\subsection{The result contract}

\begin{definition}[Process result]
A \emph{process result} $\mathcal{R}$ (schema version $1$) is the tuple
\[
   \mathcal{R} = \big(\, \mathrm{id},\, \mathrm{kind},\, \Pi,\, \mathcal{D},\, \mathcal{P},\,
   \mathcal{F},\, \mathcal{S},\, \mathcal{G},\, \Sigma \,\big),
\]
with run identifier $\mathrm{id}$; kind; parameters $\Pi$; data fingerprint $\mathcal{D}$ (directory,
dates, row/bar counts, content hash); provenance $\mathcal{P}$ (git SHA, dirty flag, timestamp); the
set of tested features $\mathcal{F}$ and skipped features $\mathcal{S}$ with reasons; a list of
findings $\mathcal{G}$; and a summary $\Sigma$. Each finding is
\[
   g = (\,\text{feature},\, \text{horizon},\, \text{metric},\, \text{value},\, \text{threshold},\,
   p,\, p_{\mathrm{adj}},\, \text{informative},\, \text{extras}\,),
\]
where $p_{\mathrm{adj}}$ is the Benjamini--Hochberg-adjusted $p$-value within the run and
\emph{informative} is the boolean verdict $\text{value} \ge \text{threshold}$ (and $p_{\mathrm{adj}}$
below the false-discovery level where applicable).
\end{definition}

The full result is written as an atomic JSON document; a queryable index row is mirrored to a local
database on a best-effort basis. Two disciplines in this contract matter for the rest of the paper.

\begin{remark}[The not-a-number partition, K2]
Before scoring, every process partitions its columns through a guard that removes dead, constant, or
too-thin columns---each skipped with a recorded reason (\code{all_nan}, \code{constant},
\code{n_valid<min})---so a missing optional feature category never crashes a run and never silently
contaminates an estimate. Information measures on a constant column are undefined, not zero; the
partition makes that explicit rather than letting a degenerate estimator return a spurious value.
\end{remark}

\begin{remark}[False-discovery scope]
\label{rem:fdr-scope}
Benjamini--Hochberg correction inside a result is \emph{per run}. The index of all runs is a
research log, \emph{not} a meta-gate: aggregating ``informative'' flags across many runs inflates
discoveries exactly as multiple testing predicts. The program-level ledger of
Section~\ref{sec:hygiene} is the honest counterpart, and the deflated-Sharpe and reality-check
gates remain authoritative for any promotion decision.
\end{remark}

\subsection{A typed-node view}

The evaluation/transform dichotomy is a convenient first cut, but it already leaks: a portfolio-book
transform performs construction rather than scoring, and a triple-barrier transform produces a
\emph{target} rather than a feature. The mature model is therefore a typed computation graph with
four node kinds---\emph{feature}, \emph{transform}, \emph{evaluation}, and \emph{target}---rather
than a two-element stack. We do not refactor preemptively; we record the target node as the first
new kind, because making the prediction target explicit (Section~\ref{sec:target}) is the single
highest-leverage change for honest information measurement.

% ============================================================
\section{Information as the Currency of Prediction}
\label{sec:info}
% ============================================================

\subsection{Why mutual information, not correlation}

Let $X$ be a feature and $Y$ a target derived from future price action. The mutual information
\[
   \MI(X;Y) \;=\; \iint p(x,y)\,\log\frac{p(x,y)}{p(x)\,p(y)}\;dx\,dy
   \;=\; H(Y) - H(Y\given X)
\]
measures \emph{any} statistical dependence, linear or not, and is invariant to monotone
reparameterization of either variable. Correlation answers ``how much, linearly''; mutual
information answers ``is there any exploitable dependence at all.'' The reason this distinction is
not cosmetic is a hard inequality.

\begin{proposition}[Information lower-bounds achievable error]
\label{prop:fano}
Let $\widehat{Y}=g(X)$ be any predictor of a discretized target $Y$ taking values in a set of size
$K$, and let $P_e=\Pr[\widehat{Y}\neq Y]$. Then Fano's inequality gives
\[
   H(Y\given X) \;\le\; H_b(P_e) + P_e\,\log(K-1),
\]
hence
\[
   P_e \;\ge\; \frac{H(Y) - \MI(X;Y) - \log 2}{\log(K-1)} .
\]
\end{proposition}

The content of Proposition~\ref{prop:fano} is that no estimator, however clever, can drive the error
below a floor set by $\MI(X;Y)$. Information is therefore not one diagnostic among several; it is the
budget. A feature with negligible mutual information cannot be rescued by a better model, and a
feature selection that maximizes information is maximizing the only quantity that bounds achievable
prediction error (Cover \& Thomas; the continuous analogue is the rate--distortion theorem).

\subsection{The cost-aware information gate}

Statistical information is necessary but not sufficient: a signal must carry enough information to
pay the round-trip trading cost. The platform encodes this as a minimum-information threshold that
every information-valued process compares against,
\[
   I_{\min} \;=\; I_{\min}\!\big(\,\text{fee},\, \sigma_r,\, \mathrm{kurt}(r)\,\big),
\]
the least mutual information a feature must share with returns for a position to clear fees and
slippage given the return scale and tail heaviness. A finding is \emph{informative} only if its
estimated information exceeds $I_{\min}$, not merely zero. This is the information-theoretic parent
of the deflated, cost-adjusted gates used downstream, and it is why a process reports bits against a
threshold rather than a bare $p$-value.

\subsection{Estimating information without manufacturing it}

Continuous mutual information is estimated with the Kraskov--St\"ogbauer--Grassberger $k$-nearest
neighbour estimator. One practical caveat is load-bearing for everything that follows.

\begin{remark}[Rank-transform before estimating]
\label{rem:ksg}
Raw-scale $k$-NN estimation manufactures a spurious information floor (empirically $\approx 0.07$
bits on planted-null data) when feature and return scales differ, because the neighbour graph is
dominated by the wider-scale coordinate. Rank/copula-transforming each input before estimation
removes the artefact and makes the estimator scale-invariant. Every information-valued process in
the framework rank-transforms its inputs; this single fix was found only because a planted-null test
(Section~\ref{sec:status}) flagged a nonzero estimate where the true value was zero.
\end{remark}

Directed dependence---which way information flows between two series---is measured by transfer
entropy $\TE(X\to Y)=\MI(Y_{t};X_{t-1}\given Y_{t-1})$ (Schreiber), with a linear (Gaussian)
default where it is tolerable and the reverse direction $\TE(Y\to X)$ retained as a control.

% ============================================================
\section{The Effective Degrees of Freedom of a Signal Book}
\label{sec:dof}
% ============================================================

The central quantity of the whole framework is not the information of any single feature but the
\emph{breadth} of the book of features that survive selection. We define it, derive it under the
natural redundancy model, and then show why its calm-regime value is a dangerous overstatement.

\subsection{The fundamental law and effective breadth}

The fundamental law of active management (Grinold; Grinold \& Kahn) relates the information ratio of
a strategy to the skill per bet and the number of independent bets,
\begin{equation}
   \IR \;\approx\; \ICbar \,\sqrt{\BR},
   \label{eq:fl}
\end{equation}
where $\ICbar$ is the average information coefficient (skill) and $\BR$ is the breadth---the number
of independent bets per period. The research programme is the pursuit of the two factors in
\eqref{eq:fl}: raise $\ICbar$ by finding real information, raise $\BR$ by combining \emph{independent}
pieces. The question is what $\BR$ equals when the pieces are correlated, which they always are.

\begin{proposition}[Effective breadth under equicorrelation]
\label{prop:neff}
Let $s_1,\dots,s_N$ be standardized signals ($\E[s_i]=0$, $\Var(s_i)=1$) with common pairwise
correlation $\corr(s_i,s_j)=\rhobar\in[0,1)$ for $i\neq j$. The equal-weight aggregate
$S=\tfrac{1}{N}\sum_{i=1}^{N}s_i$ has variance
\[
   \Var(S) \;=\; \frac{1}{N^2}\Big(\,N + N(N-1)\rhobar\,\Big)
            \;=\; \frac{1+(N-1)\rhobar}{N}.
\]
Define the \emph{effective number of independent signals} $\Neff$ as the count of mutually
independent unit-variance signals whose equal-weight aggregate has the same variance. Since $n$
independent unit-variance signals give aggregate variance $1/n$, equating $1/\Neff=\Var(S)$ yields
\begin{equation}
   \boxed{\;\Neff \;=\; \dfrac{N}{1+(N-1)\rhobar}\;}
   \label{eq:neff}
\end{equation}
and, identifying breadth with effective independence ($\BR=\Neff$) in \eqref{eq:fl},
\[
   \IR \;\approx\; \ICbar\,\sqrt{\dfrac{N}{1+(N-1)\rhobar}}.
\]
\end{proposition}

\begin{proof}
$\Var(S)=N^{-2}\big[\sum_i\Var(s_i)+\sum_{i\neq j}\cov(s_i,s_j)\big]
        =N^{-2}\big[N+N(N-1)\rhobar\big]$. Equate with the independent-case variance $1/\Neff$ and
solve.
\end{proof}

\begin{corollary}[The two limits]
As $\rhobar\to 0$, $\Neff\to N$: uncorrelated signals each count fully. As $\rhobar\to 1$,
$\Neff\to 1$: perfectly correlated signals count as one, regardless of how many were collected.
Equation~\eqref{eq:neff} is monotone decreasing in both $N$ and $\rhobar$ in the relevant regime, so
\emph{adding a redundant signal raises $N$ but cannot raise $\Neff$}---the defining property a
selection procedure must respect.
\end{corollary}

Proposition~\ref{prop:neff} is the precise sense in which this paper ``defines the degrees of
freedom mathematically.'' The raw count $N$ of significant features is a vanity metric; $\Neff$ is
the quantity that enters the achievable information ratio, and the entire orthogonalization suite of
Section~\ref{sec:ortho} exists to make $\Neff$ as close to $N$ as the data permit.

\subsection{The correlation that matters is the stress correlation}

The correlation $\rhobar$ in \eqref{eq:neff} is not a constant; it is regime-dependent,
$\rhobar=\rhobar(\text{regime})$. Microstructure signals that are nearly uncorrelated in calm
markets routinely co-move under stress, when a single liquidity factor drives every book at once.
Hence two values of \eqref{eq:neff} must be reported:
\[
   \Neff^{\text{calm}} = \frac{N}{1+(N-1)\,\rhobar_{\text{calm}}},
   \qquad
   \Neff^{\text{stress}} = \frac{N}{1+(N-1)\,\rhobar_{\text{stress}}},
   \qquad \rhobar_{\text{stress}} \ge \rhobar_{\text{calm}}.
\]

\begin{remark}[Diversification evaporates when it is needed]
A book with $\rhobar_{\text{calm}}=0.1$ but $\rhobar_{\text{stress}}=0.7$ has, for $N=10$,
$\Neff^{\text{calm}}\approx 5.3$ but $\Neff^{\text{stress}}\approx 1.4$: it offers five independent
bets in quiet markets and barely more than one exactly when a drawdown is unfolding. The gap
$\Neff^{\text{calm}}-\Neff^{\text{stress}}$ is therefore not a footnote---it \emph{is} the risk
disclosure of the book, and any honest report of breadth must state both numbers and the regime
buckets they were computed in.
\end{remark}

% ============================================================
\section{Extracting Orthogonal Information}
\label{sec:ortho}
% ============================================================

We now give the five processes that build a high-$\Neff$ book. Each is stated as a definition with
its planted-signal contract; the layer they belong to is named in Table~\ref{tab:ortho}.

\subsection{Conditional-MI greedy selection (the core extractor)}

\begin{definition}[Redundancy-aware information selection]
\label{def:cmi}
Given a feature set $F=\{f_1,\dots,f_p\}$ and target $r$, build a selected set $S_k$ greedily:
\[
   f_{(1)} = \argmax_{f\in F}\ \MI(f;r),
   \qquad
   f_{(k+1)} = \argmax_{f\in F\setminus S_k}\ \MI\big(f;\,r \given S_k\big),
\]
terminating when the marginal conditional information gain falls below
$\max\!\big(I_{\min},\,\varepsilon\big)$. Report per selected feature the gain in bits and, per
rejected runner-up, the redundancy $\MI(f;r)-\MI(f;r\given S_k)$.
\end{definition}

The conditional mutual information $\MI(f;r\given S_k)$ is the information $f$ shares with $r$ that
is \emph{not already present} in the selected set. A feature that duplicates an already-selected one
contributes $\approx 0$ and is never chosen, so the selected set is orthogonal-by-construction in
the information metric. This is strictly stronger than linear (principal-component) decorrelation,
which a nonlinear redundancy---two features that jointly but not marginally determine the target---
slips past entirely.

\begin{remark}[Relation to mRMR]
The first-order expansion of Definition~\ref{def:cmi} is the minimum-redundancy maximum-relevance
criterion, $\MI(f;r)-\tfrac{1}{|S_k|}\sum_{s\in S_k}\MI(f;s)$ (Peng, Long \& Ding; unified by Brown
et al.), and the lineage traces to Battiti's mutual-information feature selection. The conditional
form is the exact object the approximation targets.
\end{remark}

\noindent\emph{Planted contract.} On a frame containing a signal $A$, an exact noisy copy $A'$, and
an independent signal $B$, the process must select $\{$one of $A,A'$; $B\}$ and never both copies;
the cumulative-information curve must flatten after the second pick.

\subsection{Interaction information: mining synergy}

\begin{definition}[Pairwise synergy]
\label{def:synergy}
For features $X,Y$ and target $r$, the interaction information is
\[
   \mathcal{S}(X,Y;r) \;=\; \MI\big(\{X,Y\};r\big) - \MI(X;r) - \MI(Y;r).
\]
$\mathcal{S}>0$ signals \emph{synergy} (the pair jointly carries more than the sum of its parts);
$\mathcal{S}<0$ signals \emph{redundancy}.
\end{definition}

Synergy is invisible to any single-feature screen and to Definition~\ref{def:cmi}'s forward pass,
because each marginal can be near zero while the joint is large. Run over the top-$K$ features only
($K\approx 20\Rightarrow 190$ pairs, the cost of $k$-NN estimation being the binding constraint),
synergistic pairs become candidate engineered composites fed back as derived features (the partial
information decomposition of Williams \& Beer is the formal parent).

\noindent\emph{Planted contract.} On an XOR construction $r=\sgn(x)\cdot\sgn(y)$, each marginal
information is $\approx 0$ while the joint exceeds the gate; the process must flag that pair and only
that pair.

\subsection{Residualization: pure-innovation extraction}

\begin{definition}[Residualization]
\label{def:resid}
For a target feature $f$ and a conditioning set $Z$ (a market mode, a cousin feature, or a principal
component), the residual is
\[
   \mathrm{res}_f(t) = f(t) - \boldsymbol{\beta}^{\!\top} Z(t),
   \qquad
   \boldsymbol{\beta} = \argmin_{\boldsymbol{\beta}}\ \E_{\text{train}}\big[(f-\boldsymbol{\beta}^{\!\top} Z)^2\big],
\]
with $\boldsymbol{\beta}$ fit on the training prefix only. By the normal equations
$\corr(\mathrm{res}_f,Z)=0$ on the training sample, and approximately so out of sample; scoring
$\mathrm{res}_f$ against $r$ isolates what $f$ knows that $Z$ does not.
\end{definition}

Residualization is the per-feature, tradeable analogue of the set-level selection of
Definition~\ref{def:cmi}: where conditional-MI selection \emph{ranks}, residualization \emph{emits a
series}. The training-prefix fit is the no-lookahead guarantee, verified by the same holdout
perturbation test used for principal components.

\subsection{Denoised principal components and the count of real factors}

Linear orthogonalization by principal components is cheap and complementary, but on short samples
most ``components'' are noise dressed as structure. Random matrix theory says exactly how much.

\begin{proposition}[Marchenko--Pastur noise band]
\label{prop:mp}
Let $\mathbf{C}$ be the sample correlation matrix of a $T\times N$ panel of i.i.d.\ standardized
noise with aspect ratio $q=N/T\in(0,1]$. As $N,T\to\infty$ with $q$ fixed, the eigenvalues of
$\mathbf{C}$ fall in $[\lambda_-,\lambda_+]$ with
\[
   \lambda_{\pm} = \big(1 \pm \sqrt{q}\,\big)^2 .
\]
Eigenvalues exceeding $\lambda_+$ cannot be attributed to noise; the count
\[
   n_{\text{signal}} = \#\{\, i : \lambda_i > \lambda_+ \,\}
\]
estimates the number of genuine factors---the linear degrees of freedom---in the feature block.
\end{proposition}

Clipping or shrinking the eigenvalues inside $[\lambda_-,\lambda_+]$ denoises the correlation
estimate (Laloux et al.; L\'opez de Prado), and $n_{\text{signal}}$ is a direct, linear estimate of
the degrees of freedom that Proposition~\ref{prop:neff} measures in the information metric. The
principal-component transform reports $n_{\text{signal}}$ and the maximum off-diagonal holdout
correlation as its orthogonality witness.

\noindent\emph{Planted contract.} A frame of three true factors plus forty noise features must
return $n_{\text{signal}}\approx 3$, and the denoised holdout orthogonality must strictly improve on
the raw.

\subsection{The capstone: a decorrelated signal book with honest breadth}

\begin{definition}[Decorrelated signal book]
\label{def:book}
Given features that pass selection (Definition~\ref{def:cmi}), a stability gate, and a
regime-robustness gate, the signal book is built by (i) sequentially residualizing in selection
order so the surviving series are mutually orthogonal, (ii) walk-forward $z$-scoring them into
$\code{sig_*}$ columns, and (iii) reporting, as the book's contract,
\[
   \Neff = \frac{N}{1+(N-1)\rhobar}, \qquad
   \IR_{\text{expected}} = \ICbar\,\sqrt{\Neff \cdot \text{bars/year}},
\]
computed at \emph{both} calm and stress $\rhobar$, together with the full provenance chain of every
upstream run that admitted each entry.
\end{definition}

The book is the object an algorithm consumes: an algorithm whose required inputs are the
$\code{sig_*}$ columns and whose step is risk-parity weighting. The division of labour is the thesis
of the framework in one line: \emph{processes certify what is combinable; algorithms combine.}

\begin{table}[h]
\centering
\renewcommand{\arraystretch}{1.3}
\setlength{\tabcolsep}{5pt}
\begin{tabular}{>{\raggedright\arraybackslash}p{4.0cm}
                >{\raggedright\arraybackslash}p{2.1cm}
                >{\raggedright\arraybackslash}p{8.1cm}}
\hline
\textbf{Process} & \textbf{Kind} & \textbf{Role in raising $\Neff$} \\
\hline
\code{cmi_select}          & evaluation & selects an information-orthogonal basis at the source \\
\code{interaction_synergy} & evaluation & recovers joint information no marginal screen sees \\
\code{lead_lag_te}         & evaluation & cross-symbol directed edges---the cheapest decorrelation \\
\code{residualize}         & transform  & emits the pure innovation of a feature beyond a set \\
\code{pca_combo}\,$+$\,MP  & transform  & linear orthogonal composites; counts real factors \\
\code{signal_book}         & transform  & assembles the mutually orthogonal book; reports $\Neff$ \\
\hline
\end{tabular}
\caption{The orthogonalization suite. Each process either raises the effective breadth $\Neff$ or
certifies it. \code{pca_combo} is implemented; the others are specified (Section~\ref{sec:status}).}
\label{tab:ortho}
\end{table}

% ============================================================
\section{Defining the Target: What ``Price Action'' Means}
\label{sec:target}
% ============================================================

Mutual information is always information \emph{about something}. A pervasive but silent error is to
let every process default to simple forward returns while claiming to measure ``information about
price action.'' Price action is not one thing, and the literature has studied features against many
distinct targets. We make the target a first-class node.

\begin{definition}[Target node]
A \emph{target} is a map $\tau:(\text{prices},\,\text{bars},\,\text{horizon})\mapsto \R^{T}$ with a
declared kind (continuous, binary, categorical), resolved from a registry and carried in the process
context, so that every process computes $\MI(f;\tau)$ against one explicit, named target rather than
an implicit default. Closing this gap is the highest-leverage single change in the framework: it
makes ``mutual information between a feature and price action'' honest about \emph{which} price
action.
\end{definition}

Table~\ref{tab:targets} is the resulting design surface: a run is the quadruple
(process $\times$ target $\times$ horizon $\times$ symbol).

\begin{table}[h]
\centering
\renewcommand{\arraystretch}{1.3}
\setlength{\tabcolsep}{6pt}
\begin{tabular}{>{\raggedright\arraybackslash}p{4.3cm}
                >{\raggedright\arraybackslash}p{4.6cm}
                >{\raggedright\arraybackslash}p{4.9cm}}
\hline
\textbf{Aspect} & \textbf{Definition} & \textbf{Literature / status} \\
\hline
Direction        & $\sgn(r_h)$                                  & classification; partial support \\
Magnitude        & $r_h$, $|r_h|$                               & rate--distortion; \emph{shipped default} \\
Volatility       & future realized variance over $h$            & HAR (Corsi); a target, not a scaler \\
Path             & triple-barrier label, time-to-touch          & L\'opez de Prado; partial \\
Regime change    & $\Pr(\text{switch})$                         & Hamilton; unbuilt \\
Jump / tail      & Lee--Mykland statistic                       & jump literature; algorithm only \\
Fill-conditional & move at the transacted price                 & Kyle/Stoikov; the missing bridge \\
\hline
\end{tabular}
\caption{An aspect$\times$horizon taxonomy of price-action targets (horizons: tick--second,
minute--bar, multi-day). Two cells are both unbuilt and high value: \emph{volatility-as-target}
(a whole tradeable family the platform cannot currently screen for) and \emph{fill-conditional move}
---the target defined on the price one actually transacts at, where the mid-price edge of
$\approx 0.45$ is known to collapse to $\approx 0.03$. The latter is the execution-realism gap made
measurable.}
\label{tab:targets}
\end{table}

% ============================================================
\section{Minting Feature Degrees of Freedom}
\label{sec:ops}
% ============================================================

Orthogonalization raises $\Neff$ toward the number of distinct bets present in the feature set; it
cannot create a bet that no feature expresses. The complementary lever is a transform process,
$\code{feature_ops}$, that applies operators to a chosen base-column set and emits new
$\code{op_*}$ columns---candidate new degrees of freedom---each chainable into the evaluation and
selection processes above. Only the operators that are genuinely missing from the existing feature
set are proposed; the high-value, not-yet-built ones are:

\begin{enumerate}
  \item \textbf{Fractional differentiation} $(1-B)^{d}$, $0<d<1$ (L\'opez de Prado, after the long-memory
        literature). Integer differencing destroys the memory that carries predictive content;
        fractional differencing makes a level series stationary while \emph{preserving} memory,
        so downstream information and IC estimates are both honest (stationary) and not amputated.
        This is the single highest-value missing operator.
  \item \textbf{Leaky / cumulative integration} $y_t = \varrho\,y_{t-1} + x_t$, accumulating a
        memoryless flow or imbalance into a persistent pressure state for longer-horizon targets.
  \item \textbf{Causal spectral features}: rolling band-power, dominant frequency, and rolling
        spectral entropy, exposing oscillatory structure as a persisted feature rather than as a
        one-off analysis.
  \item \textbf{Distribution-free normalization}: rolling rank/quantile and robust (median/MAD)
        $z$-scores, which both help the rank-based information estimator of
        Remark~\ref{rem:ksg} and are fat-tail robust.
  \item \textbf{Nonlinear bases}: polynomial/spline expansions and a sign$\times$magnitude split, so
        that even linear processes catch curvature and direction-versus-size asymmetry.
\end{enumerate}

\begin{remark}[Operators sharpen, rather than relax, the multiple-testing problem]
\label{rem:ledger-ops}
Each operator multiplies the column space: one base feature becomes many $\code{op_*}$ candidates.
A $\code{feature_ops}$ run must therefore log how many columns it minted, so the program-level
false-discovery ledger of Section~\ref{sec:hygiene} can count them. Minting degrees of freedom and
accounting for them are the same obligation.
\end{remark}

% ============================================================
\section{Estimation Hygiene: Keeping the Breadth Estimate Honest}
\label{sec:hygiene}
% ============================================================

A measured information coefficient or mutual information is an estimate, and estimates of small
quantities on short, overlapping, heavy-tailed samples are biased upward by selection. Three
processes protect the degrees-of-freedom estimate from being fiction.

\paragraph{Permutation null (reality check).} For a candidate set or a prior run, rebuild the null
by stationary block bootstrap of returns---which preserves autocorrelation and the overlap structure
that naive shuffling destroys (Politis \& Romano)---over $M$ resamples, and recompute the headline
metric. Report, per feature, $p_{\text{perm}}=M^{-1}\#\{\text{metric}_{\text{null}}\ge
\text{metric}_{\text{obs}}\}$, and at the set level the \emph{excess discoveries}
$n_{\text{obs}}-\E[n_{\text{null}}]$ (White's reality check). A set whose excess is $\approx 0$ is
indistinguishable from luck regardless of any single feature's $p$-value.

\paragraph{Stability and decay.} Signals are perishable inventory. A non-stationarity audit computes
a calendar-windowed (not expanding) IC series, detects breaks by a CUSUM/Page--Hinkley test, and
measures sign-consistency and the half-life of $|\IC|$ with window age. A feature is admissible only
if no break is detected and its sign is consistent across sub-periods.

\paragraph{Program-level ledger.} Per Remark~\ref{rem:fdr-scope}, the number that decides whether
discoveries are real is the \emph{total} count of hypotheses the platform has ever tested. A ledger
the runner appends to---(process, target, columns minted, timestamp, git SHA)---is the meta-gate at
program scope, the generalization of the reality check from one run to the whole research history.
Until it exists, every information number is un-budgeted against the platform's cumulative search
(Bailey, Borwein, L\'opez de Prado \& Zhu).

% ============================================================
\section{Verification and Implementation Status}
\label{sec:status}
% ============================================================

\subsection{The planted-signal discipline}

No process touches real data before it passes a \emph{planted-signal} conformance test: a synthetic
frame with a known signal column and known noise columns, on which the process must flag the signal
and must \emph{not} flag the noise. This is not optional hygiene---three estimator bugs (including
the scale artefact of Remark~\ref{rem:ksg}) were caught only by planted tests, each of which would
otherwise have manufactured a discovery. A transform additionally passes a no-lookahead test (a
holdout perturbation must not change in-sample outputs). A real-parquet smoke test on the latest day,
with dead columns skipped and reasons recorded, is the final pre-commit gate.

\subsection{What is built and what is specified}

In the interest of honest reporting, Table~\ref{tab:status} separates the implemented contract and
processes from the specified-but-unbuilt mechanisms. The framework, its result schema, its
persistence, and seven processes are implemented; the orthogonalization suite that this paper
centres on is specified with its planted contracts and is the immediate build queue.

\begin{table}[h]
\centering
\renewcommand{\arraystretch}{1.3}
\setlength{\tabcolsep}{6pt}
\begin{tabular}{>{\raggedright\arraybackslash}p{6.4cm}
                >{\raggedright\arraybackslash}p{8.0cm}}
\hline
\textbf{Implemented} & \textbf{Specified (planted contract written, not yet built)} \\
\hline
process contract + result schema + persistence    & \code{cmi_select} (conditional-MI selection) \\
\code{ic_horizon} (expanding Spearman IC sweep)    & \code{interaction_synergy} (synergy mining) \\
\code{mi_ksg} (KSG mutual information + CMI)        & \code{lead_lag_te} (cross-symbol TE) \\
\code{transfer_entropy} (directed information)      & \code{residualize} (pure innovation) \\
\code{spectral} (PSD / persistence)                & \code{permutation_null}, \code{stability_decay} \\
\code{ml_importance} (walk-forward importance)      & \code{regime_conditional}, \code{signal_book} \\
\code{triple_barrier}, \code{pca_combo}            & \code{targets} registry, \code{feature_ops} operators \\
\hline
\end{tabular}
\caption{Implementation status. This is a methods/framework paper: the orthogonalization and
degrees-of-freedom mathematics are the contribution, the implemented processes demonstrate the
contract, and empirical estimation of $\Neff$ on a clean multi-day window is specified as future
work---meaningful breadth estimates require more independent observations than a thin sample
provides.}
\label{tab:status}
\end{table}

\subsection{The end-to-end pipeline}

The processes assemble into a single composition, each arrow a runner chain:
\[
\begin{array}{c}
   \underbrace{\code{feature_ops}}_{\text{mint DOF}}
   \;\to\;
   \underbrace{\code{cmi_select}}_{\text{pick orthogonal}}
   \;\to\;
   \underbrace{\text{Target}}_{\text{which price action}}
   \;\to\;
   \\[2.4ex]
   \underbrace{\text{evaluation}}_{\text{score in bits / IC}}
   \;\to\;
   \underbrace{\code{signal_book}}_{\text{book}+\Neff}
   \;\to\;
   \underbrace{\text{algorithm}}_{\text{risk-parity combine}} .
\end{array}
\]
Two maturity gaps remain explicitly open. The first is execution realism: there is no fill/cost
model node, and the fill-conditional target of Table~\ref{tab:targets} is the hook for the missing
fourth citizen, where a mid-price edge becomes a transacted-price edge or does not. The second is the
program-level multiple-testing ledger of Section~\ref{sec:hygiene}. Both are named here precisely so
that no estimate downstream is mistaken for more than it is.

% ============================================================
\section{Conclusion}
\label{sec:conclusion}
% ============================================================

We have argued that the unit of microstructure research should be the \emph{process}: a typed,
registry-bound analytical contract that certifies whether features carry non-redundant information
about an explicit target. Around that unit we developed the two pieces of mathematics it exists to
serve. Mutual information, justified by Fano's inequality as a hard floor on achievable error, is
the right currency; and the effective degrees of freedom $\Neff=N/[1+(N-1)\rhobar]$, derived from
the fundamental law of active management, is the right measure of a book's worth---with the
stress-regime correlation, not the calm one, as the binding input. Orthogonal information is then
extracted by five composable processes (conditional-MI selection, synergy mining, residualization,
denoised principal components, and a capstone signal book), the target is made explicit along an
aspect$\times$horizon taxonomy, and new degrees of freedom are minted by a derivative-operator layer
under a program-level false-discovery ledger. The result is a single contract under which
``does this feature add anything?'' and ``how many independent bets do we really have?'' become
routine, testable, reproducible queries rather than one-off scripts. The empirical estimation of
$\Neff$ on a clean longitudinal window, and the execution-realism node that the fill-conditional
target anticipates, are the natural next steps---and, by construction, each will arrive behind its
own planted-signal gate.

% ============================================================
\section*{References}
\addcontentsline{toc}{section}{References}
% ============================================================

\begin{enumerate}

  \item Bailey, D.H., Borwein, J.M., L\'opez de Prado, M.\ \& Zhu, Q.J.\ (2014). Pseudo-mathematics
        and financial charlatanism: the effects of backtest overfitting on out-of-sample
        performance. \textit{Notices of the AMS}, 61(5), 458--471.

  \item Battiti, R.\ (1994). Using mutual information for selecting features in supervised neural net
        learning. \textit{IEEE Transactions on Neural Networks}, 5(4), 537--550.

  \item Benjamini, Y.\ \& Hochberg, Y.\ (1995). Controlling the false discovery rate: a practical and
        powerful approach to multiple testing. \textit{Journal of the Royal Statistical Society,
        Series B}, 57(1), 289--300.

  \item Brown, G., Pocock, A., Zhao, M.-J.\ \& Luj\'an, M.\ (2012). Conditional likelihood
        maximisation: a unifying framework for information theoretic feature selection.
        \textit{Journal of Machine Learning Research}, 13, 27--66.

  \item Cont, R.\ (2001). Empirical properties of asset returns: stylized facts and statistical
        issues. \textit{Quantitative Finance}, 1(2), 223--236.

  \item Corsi, F.\ (2009). A simple approximate long-memory model of realized volatility.
        \textit{Journal of Financial Econometrics}, 7(2), 174--196.

  \item Cover, T.M.\ \& Thomas, J.A.\ (2006). \textit{Elements of Information Theory}, 2nd ed.\
        Wiley-Interscience.

  \item Grinold, R.C.\ (1989). The fundamental law of active management. \textit{Journal of Portfolio
        Management}, 15(3), 30--37.

  \item Grinold, R.C.\ \& Kahn, R.N.\ (2000). \textit{Active Portfolio Management}, 2nd ed.\
        McGraw-Hill.

  \item Hamilton, J.D.\ (1989). A new approach to the economic analysis of nonstationary time series
        and the business cycle. \textit{Econometrica}, 57(2), 357--384.

  \item Hasbrouck, J.\ (1995). One security, many markets: determining the contributions to price
        discovery. \textit{Journal of Finance}, 50(4), 1175--1199.

  \item Kraskov, A., St\"ogbauer, H.\ \& Grassberger, P.\ (2004). Estimating mutual information.
        \textit{Physical Review E}, 69(6), 066138.

  \item Kyle, A.S.\ (1985). Continuous auctions and insider trading. \textit{Econometrica}, 53(6),
        1315--1335.

  \item Laloux, L., Cizeau, P., Bouchaud, J.-P.\ \& Potters, M.\ (1999). Noise dressing of financial
        correlation matrices. \textit{Physical Review Letters}, 83(7), 1467--1470.

  \item Lee, S.S.\ \& Mykland, P.A.\ (2008). Jumps in financial markets: a new nonparametric test and
        jump dynamics. \textit{Review of Financial Studies}, 21(6), 2535--2563.

  \item L\'opez de Prado, M.\ (2018). \textit{Advances in Financial Machine Learning}. Wiley.
        (Fractional differentiation, Ch.\ 5; triple-barrier labelling, Ch.\ 3; correlation-matrix
        denoising.)

  \item Marchenko, V.A.\ \& Pastur, L.A.\ (1967). Distribution of eigenvalues for some sets of random
        matrices. \textit{Mathematics of the USSR-Sbornik}, 1(4), 457--483.

  \item Peng, H., Long, F.\ \& Ding, C.\ (2005). Feature selection based on mutual information:
        criteria of max-dependency, max-relevance, and min-redundancy. \textit{IEEE Transactions on
        Pattern Analysis and Machine Intelligence}, 27(8), 1226--1238.

  \item Politis, D.N.\ \& Romano, J.P.\ (1994). The stationary bootstrap. \textit{Journal of the
        American Statistical Association}, 89(428), 1303--1313.

  \item Schreiber, T.\ (2000). Measuring information transfer. \textit{Physical Review Letters},
        85(2), 461--464.

  \item Shannon, C.E.\ (1948). A mathematical theory of communication. \textit{Bell System Technical
        Journal}, 27, 379--423, 623--656.

  \item White, H.\ (2000). A reality check for data snooping. \textit{Econometrica}, 68(5),
        1097--1126.

  \item Williams, P.L.\ \& Beer, R.D.\ (2010). Nonnegative decomposition of multivariate information.
        \textit{arXiv:1004.2515}.

\end{enumerate}

\end{document}
